\chapter{Výsledky a~diskuse}
\label{chap:results}

Tato kapitola shrnuje výsledky jednotlivých fází výzkumu a~diskutuje jejich implikace pro marketingovou strategii značky Berrie při vstupu na český trh.

% ============================================================================
\section{Vyhodnocení skupinových diskusí}
\label{sec:res_fgd}

\textit{Tato sekce bude doplněna po realizaci a~analýze skupinových diskusí.}

\subsection{Vnímání kategorie a~první dojem}

% TODO: Spontánní asociace s kategorií mraženého ovoce v čokoládě
% TODO: Povědomí o Franui a odlišení od tradičních mražených produktů
% TODO: Emocionální reakce na koncept produktu

\subsection{Senzorické hodnocení}

% TODO: Hodnocení textury v čase (5, 10, 15 min)
% TODO: Chuťová rovnováha čokoláda vs. ovoce
% TODO: Preference variant (borůvky, jahody, banán, mix)

\subsection{Positioning a~vnímání značky}

% TODO: Reakce na vizuální stránku obalu
% TODO: Důvěra v lokální původ a BIO kvalitu
% TODO: Vnímání ceny

\subsection{Rozdíly mezi segmenty}

% TODO: Porovnání výsledků Gen Z vs. Matky
% Očekávané rozdíly:
% - Gen Z: důraz na vizuální stránku, sociální sítě, prémiovost
% - Matky: důraz na složení, zdravotní aspekty, vhodnost pro děti

% ============================================================================
\section{Vyhodnocení dotazníkového šetření}
\label{sec:res_survey}

\textit{Tato sekce bude doplněna po sběru a~analýze kvantitativních dat.}

\subsection{Demografický profil vzorku}

% TODO: Tabulka s demografickým rozložením respondentů
% (věk, pohlaví, region, příjem)

\subsection{Chuťové preference}

% TODO: Které druhy ovoce jsou preferovány?
% TODO: Grafy preferencí (koláčový/sloupcový)

\subsection{Nákupní chování}

% TODO: Příležitosti ke koupi
% TODO: Distribuční preference (online vs. supermarket)
% TODO: Frekvence potenciálního nákupu

% ============================================================================
\section{Cenová analýza}
\label{sec:res_pricing}

\textit{Tato sekce bude doplněna po vyhodnocení cenové části dotazníku.}

\subsection{Výsledky Price Sensitivity Meter}

% TODO: Graf kumulativních distribučních křivek (Van Westendorp)
% TODO: Identifikace:
%   - Optimální cenový bod (OPP)
%   - Bod cenové indiference (IDP)
%   - Přijatelné cenové rozmezí
% TODO: Srovnání s cenou Franui (~160 Kč / 150g)

\subsection{Cenová doporučení}

% TODO: Doporučená cenová strategie na základě PSM výsledků

% ============================================================================
\section{Výsledky senzorické analýzy}
\label{sec:res_sensory}

\textit{Tato sekce bude doplněna po vyhodnocení Central Location Testu.}

% TODO: Hodnocení textury v časových intervalech
% TODO: Srovnání Berrie vs. Franui (pokud bylo realizováno)
% TODO: Celkový senzorický profil produktu

% ============================================================================
\section{Implikace pro marketingovou strategii}
\label{sec:res_implications}

Na základě dosavadních zjištění z~analýzy trhu a~konkurenčního prostředí lze formulovat předběžná doporučení pro marketingovou strategii značky Berrie. Tato doporučení budou zpřesněna po dokončení primárního výzkumu.

\subsection{Produktová strategie}

Značka Berrie by měla vstoupit na trh s~portfoliem minimálně tří variant (borůvky, jahody, banán), čímž se odliší od Franui, které nabízí pouze maliny. Důraz na šíři portfolia umožňuje oslovit různé chuťové preference a~zvyšuje šanci na opakovaný nákup.

\subsection{Komunikační strategie}

S~ohledem na úspěch Franui na sociálních sítích je doporučeno:
\begin{itemize}
  \item vybudování silné přítomnosti na TikToku a~Instagramu s~vizuálně atraktivním obsahem,
  \item spolupráce s~food a~lifestyle influencery v~české komunitě,
  \item zdůraznění příběhu lokálního původu jako diferenciačního faktoru,
  \item sampling a~ochutnávkové akce v~místě prodeje.
\end{itemize}

\subsection{Distribuční strategie}

Vstup prostřednictvím platformy Rohlík.cz je doporučen jako primární distribuční kanál, neboť:
\begin{itemize}
  \item platforma je silným kanálem pro uvedení prémiových novinek,
  \item cílová skupina (urbánní, zdraví vědomí spotřebitelé) se výrazně překrývá s~uživatelskou základnou Rohlíku,
  \item online distribuce umožňuje nižší vstupní náklady oproti okamžitému vstupu do kamenných řetězců.
\end{itemize}

% ============================================================================
\section{Ekonomický náhled}
\label{sec:res_economics}

Rozpočet výzkumného projektu je stanoven na minimálně 90\,000~Kč. Struktura nákladů je uvedena v~Tabulce~\ref{tab:budget}.

\begin{table}[H]
\centering
\caption{Rozpočet marketingového výzkumu}
\label{tab:budget}
\begin{tabular}{lr}
\toprule
\textbf{Položka} & \textbf{Částka (Kč)} \\
\midrule
Online dotazníkové šetření (CAWI) & 42\,000 \\
Skupinové diskuse (pronájem, moderátor, odměny) & 22\,000 \\
Senzorické testování (CLT) & 16\,000 \\
Suroviny a vzorky (Berrie + Franui) & 10\,000 \\
\midrule
\textbf{Celkem} & \textbf{90\,000} \\
\bottomrule
\end{tabular}
\end{table}

Tato investice má za cíl snížit riziko neúspěchu při vstupu na trh. Výzkum umožní optimalizovat cenovou strategii, komunikační sdělení a~produktové portfolio ještě před spuštěním komerční výroby a~distribuce.

% ============================================================================
\section{Limity výzkumu}
\label{sec:res_limits}

Při interpretaci výsledků je nutné zohlednit následující omezení:
\begin{itemize}
  \item Kvalitativní výzkum neposkytuje statisticky zobecnitelné výsledky; zjištění z~focus groups slouží primárně jako vstup pro kvantitativní fázi.
  \item Kvótní výběr respondentů v~kvantitativní fázi nemusí plně odrážet strukturu cílové populace.
  \item Výběrové kritérium (aktivní uživatelé sociálních sítí) může systematicky vylučovat spotřebitele, kteří nepoužívají Instagram či TikTok, ale mohli by být potenciálními zákazníky.
  \item Senzorické testování probíhá za standardizovaných podmínek \gls{clt}, které se mohou lišit od reálné konzumace v~domácím prostředí.
  \item Výzkum se zaměřuje na český trh; výsledky nelze bez dalšího ověření přenášet na jiné trhy.
\end{itemize}
